\begin{thebibliography}{99}
\bibitem[1]{published-1}
S. Aaronson and A. Wigderson.


Algebrization: A new barrier in complexity theory.


ACM Transactions on Computation Theory (TOCT), 1(1):1–54,
2009.


doi:10.1145/1490270.1490272.

\bibitem[2]{published-2}
A. Abdesselam and J. Chipalkatti.


Brill–Gordan loci, transvectants and an analogue of the Foulkes
conjecture.


Advances in Mathematics, 208(2):491–520, 2007.


doi:10.1016/j.aim.2006.03.003.

\bibitem[3]{published-3}
A. Abdesselam, C. Ikenmeyer, and G. Royle.


16,051 formulas for Ottaviani’s invariant of cubic threefolds.


J. Algebra, 447:649–663, 2016.


doi:10.1016/j.jalgebra.2015.11.009.

\bibitem[4]{published-4}
G. E. Baird and L. C. Biedenharn.


On the representations of the semisimple Lie groups. II.


J. Mathematical Phys., 4:1449–1466, 1963.


doi:10.1063/1.1703926.

\bibitem[5]{published-5}
T. Baker, J. Gill, and R. Solovay.


Relativizations of the $\mathcal{P}=?\mathcal{NP}$ question.


SIAM Journal on computing, 4(4):431–442, 1975.


doi:10.1137/0204037.

\bibitem[6]{published-6}
F. Berdjis.


A criterion for completeness of Casimir operators.


J. Math. Phys., 22:1851–1856, 1981.


doi:10.1063/1.525156.

\bibitem[7]{published-7}
D. Bini, M. Capovani, G. Lotti, and F. Romani.


$O(n^{2.7799})$ complexity for $n\times{}n$ approximate matrix
multiplication.


Inform. Process. Lett., 8(5):234–235, 1979.


doi:10.1016/0020-0190(79)90113-3.

\bibitem[8]{published-8}
M. Bläser, M. Christandl, and J. Zuiddam.


The border support rank of two-by-two matrix multiplication is seven.


Chicago Journal OF Theoretical Computer Science, 5:1–16, 2018.

\bibitem[9]{published-9}
M. Bläser, J. Dörfler, and C. Ikenmeyer.


On the Complexity of Evaluating Highest Weight Vectors.


36th Computational Complexity Conference (CCC 2021),
200:29:1–29:36, 2021.


doi:10.4230/LIPIcs.CCC.2021.29.

\bibitem[10]{published-10}
N. Bourbaki.


Lie groups and Lie algebras. Chapters 7–9.


Elements of Mathematics (Berlin). Springer-Verlag, Berlin, 2005.


Translated from the 1975 and 1982 French originals by Andrew
Pressley.

\bibitem[11]{published-11}
P. Breiding, R. Hodges, C. Ikenmeyer, and M. Michałek.


Equations for GL invariant families of polynomials.


Vietnam Journal of Mathematics, 50(2):545–556, 2022.


doi:10.1007/s10013-022-00549-4.

\bibitem[12]{published-12}
P. Bürgisser.


Completeness and reduction in algebraic complexity theory,
volume 7.


Springer Science \& Business Media, 2000.


doi:10.1007/978-3-662-04179-6.

\bibitem[13]{published-13}
P. Bürgisser.


Completeness classes in algebraic complexity theory.


arXiv:2406.06217, 2024.

\bibitem[14]{published-14}
P. Bürgisser, M. Clausen, and M. A. Shokrollahi.


Algebraic complexity theory, volume 315.


Springer Science \& Business Media, 2013.

\bibitem[15]{published-15}
P. Bürgisser and C. Ikenmeyer.


Geometric complexity theory and tensor rank.


In Proceedings of the 43rd Annual ACM Symp. on Th. of Comp.,
pages 509–518, New York, 2011. STOC ’11, ACM.


doi:10.1145/1993636.1993704.

\bibitem[16]{published-16}
P. Bürgisser and C. Ikenmeyer.


Explicit Lower Bounds via Geometric Complexity Theory.


In Proceedings of the 45th Annual ACM Symp. on Th. of Comp.,
pages 141–150. ACM, 2013.


doi:10.1145/2488608.2488627.

\bibitem[17]{published-17}
P. Bürgisser and C. Ikenmeyer.


Fundamental invariants of orbit closures.


J. Algebra, 477:390–434, 2017.


doi:10.1016/j.jalgebra.2016.12.035.

\bibitem[18]{published-18}
P. Bürgisser, C. Ikenmeyer, and G. Panova.


No occurrence obstructions in geometric complexity theory.


J. Amer. Math. Soc., 32(1):163–193, 2019.


doi:10.1090/jams/908.

\bibitem[19]{published-19}
P. Bürgisser, J. M. Landsberg, L. Manivel, and J. Weyman.


An overview of mathematical issues arising in the Geometric
Complexity Theory approach to $VP\neq VNP$.


SIAM J. Comput., 40(4):1179–1209, 2011.


doi:10.1137/090765328.

\bibitem[20]{published-20}
A. Cayley.


On the theory of linear transformations.


Cambridge Math. J., iv:193–209, 1845.

\bibitem[21]{published-21}
P. Chatterjee, M. Kumar, C. Ramya, R. Saptharishi, and A. Tengse.


On the Existence of Algebraically Natural Proofs.


In 2020 IEEE 61st Annual Symp. on Found. of Comp. Sc. (FOCS),
pages 870–880, 2020.


Full version at arXiv:2004.14147.


doi:10.1109/FOCS46700.2020.00085.

\bibitem[22]{published-22}
X. Chen, N. Kayal, and A. Wigderson.


Partial derivatives in arithmetic complexity and beyond.


Found. Trends Theor. Comput. Sci., 6(1–2):1–138, 2011.


doi:10.1561/0400000043.

\bibitem[23]{published-23}
M.-W. Cheung, C. Ikenmeyer, and S. Mkrtchyan.


Symmetrizing tableaux and the 5th case of the Foulkes conjecture.


J. Symbolic Comp., 80:833–843, 2017.


doi:10.1016/j.jsc.2016.09.002.

\bibitem[24]{published-24}
J. Dörfler, C. Ikenmeyer, and G. Panova.


On geometric complexity theory: Multiplicity obstructions are
stronger than occurrence obstructions.


SIAM Journal on Applied Algebra and Geometry, 4(2):354–376,
2020.


doi:10.4230/LIPIcs.ICALP.2019.51.

\bibitem[25]{published-25}
P. Dutta, P. Dwivedi, and N. Saxena.


Demystifying the border of depth-3 algebraic circuits.


In 2021 IEEE 62nd Annual Symposium on Foundations of Computer
Science (FOCS), pages 92–103. IEEE, 2021.


doi:10.1109/FOCS52979.2021.00018.

\bibitem[26]{published-26}
K. Efremenko, A. Garg, R. Oliveira, and A. Wigderson.


Barriers for Rank Methods in Arithmetic Complexity.


In 9th Innovations in Theoretical Computer Science Conference
(ITCS 2018), volume 94 of Leibniz International Proceedings in
Informatics (LIPIcs), pages 1:1–1:19, Dagstuhl, Germany, 2018. Schloss
Dagstuhl – Leibniz-Zentrum für Informatik.


doi:10.4230/LIPIcs.ITCS.2018.1.

\bibitem[27]{published-27}
M. A. Forbes.


Polynomial identity testing of read-once oblivious algebraic
branching programs.


PhD thesis, Massachusetts Institute of Technology, 2014.

\bibitem[28]{published-28}
M. A. Forbes.


Deterministic divisibility testing via shifted partial derivatives.


In 2015 IEEE 56th Annual Symposium on Foundations of Computer
Science, pages 451–465. IEEE, 2015.


doi:10.1109/FOCS.2015.35.

\bibitem[29]{published-29}
M. A. Forbes.


Some concrete questions on the border complexity of polynomials.


Presentation given at the Workshop on Algebraic Complexity Theory
WACT 2016 in Tel Aviv, 2016.


URL: https://www.youtube.com/watch?v=1HMogQIHT6Q.

\bibitem[30]{published-30}
M. A. Forbes.


Low-Depth Algebraic Circuit Lower Bounds over Any Field.


In 39th Computational Complexity Conference, CCC 2024, July
22-25, 2024, volume 300 of LIPIcs, pages 31:1–31:16. Schloss
Dagstuhl – Leibniz-Zentrum für Informatik, 2024.


doi:10.4230/LIPICS.CCC.2024.31.

\bibitem[31]{published-31}
M. A. Forbes, A. Shpilka, and B. L. Volk.


Succinct hitting sets and barriers to proving lower bounds for
algebraic circuits.


Theory of Computing, 14(1):1–45, 2018.


doi:10.4086/toc.2018.v014a018.

\bibitem[32]{published-32}
W. Fulton and J. Harris.


Representation theory: a first course, volume 129 of Graduate Texts in Mathematics.


Springer-Verlag, New York, 1991.

\bibitem[33]{published-33}
M. Gałązka.


Vector bundles give equations of cactus varieties.


Lin. Alg. Appl., 521:254–262, 2017.


doi:10.1016/j.laa.2016.12.005.

\bibitem[34]{published-34}
A. Garg, C. Ikenmeyer, V. Makam, R. Oliveira, M. Walter, and A. Wigderson.


Search problems in algebraic complexity, GCT, and hardness of
generators for invariant rings.


In 35th Computational Complexity Conference, page 1, 2020.

\bibitem[35]{published-35}
I. M. Gelfand and M. L. Tsetlin.


Finite-dimensional representations of the group of unimodular
matrices.


Dokl. Akad. Nauk SSSR, 71(8):825, 1950.

\bibitem[36]{published-36}
F. Gesmundo, P. Ghosal, C. Ikenmeyer, and V. Lysikov.


Degree-Restricted Strength Decompositions and Algebraic Branching
Programs.


In 42nd IARCS Ann. Conf. Found. Soft. Tech. and TCS (FSTTCS
2022), volume 250 of Leibniz International Proceedings in Informatics
(LIPIcs), pages 20:1–20:15, Dagstuhl, Germany, 2022. Schloss Dagstuhl –
Leibniz-Zentrum für Informatik.


doi:10.4230/LIPIcs.FSTTCS.2022.20.

\bibitem[37]{published-37}
F. Gesmundo, Y.-I. Han, and B. Lovitz.


Linear preservers of secant varieties and other varieties of
tensors.


J. Symbolic Comp., 131:102449, 2025.


doi:10.1016/j.jsc.2025.102449.

\bibitem[38]{published-38}
F. Gesmundo, C. Ikenmeyer, and G. Panova.


Geometric complexity theory and matrix powering.


Diff. Geom. Appl., 55:106–127, 2017.


doi:10.1016/j.difgeo.2017.07.001.

\bibitem[39]{published-39}
J. A. Grochow.


Unifying known lower bounds via geometric complexity theory.


computational complexity, 24:393–475, 2015.


doi:10.1007/s00037-015-0103-x.

\bibitem[40]{published-40}
J. A. Grochow, M. Kumar, M. Saks, and S. Saraf.


Towards an algebraic natural proofs barrier via polynomial identity
testing.


arXiv, 2017.


arXiv:1701.01717.

\bibitem[41]{published-41}
A. Gupta, P. Kamath, N. Kayal, and R. Saptharishi.


Arithmetic circuits: A chasm at depth three.


Electronic Colloquium on Computational Complexity (ECCC),
20:26, 2013.


doi:10.1109/FOCS.2013.68.

\bibitem[42]{published-42}
B. C. Hall.


Lie Groups, Lie Algebras, and Representations: An
Elementary Introduction, volume 222 of Graduate Texts in
Mathematics.


Springer, New York, 2010.

\bibitem[43]{published-43}
J. Hauenstein, C. Ikenmeyer, and J. M. Landsberg.


Equations for lower bounds on border rank.


Experimental Mathematics, 22(4):372–383, 2013.


doi:10.1080/10586458.2013.825892.

\bibitem[44]{published-44}
J. E. Humphreys.


Introduction to Lie algebras and representation theory,
volume 9 of Graduate Texts in Mathematics.


Springer-Verlag, New York, 1978.

\bibitem[45]{published-45}
J. E. Humphreys.


Linear algebraic groups, volume 21.


Springer Science \& Business Media, 2012.

\bibitem[46]{published-46}
C. Ikenmeyer and U. Kandasamy.


Implementing geometric complexity theory: On the separation of orbit
closures via symmetries.


In Proceedings of the 52nd Annual ACM SIGACT Symposium on
Theory of Computing, pages 713–726, 2020.


doi:10.1145/3357713.3384257.

\bibitem[47]{published-47}
C. Ikenmeyer and G. Panova.


Rectangular Kronecker coefficients and plethysms in geometric
complexity theory.


Adv. Math., 319:40–66, 2017.


doi:10.1016/j.aim.2017.08.024.

\bibitem[48]{published-48}
C. Ikenmeyer and A. Sanyal.


A note on VNP-completeness and border complexity.


Information Processing Letters, 176:106243, 2022.


doi:10.1016/j.ipl.2021.106243.

\bibitem[49]{published-49}
A. G. Khovanskiĭ.


Fewnomials, volume 88.


American Mathematical Soc., 1991.

\bibitem[50]{published-50}
M. Kumar.


A quadratic lower bound for homogeneous algebraic branching
programs.


computational complexity, 28:409–435, 2019.


doi:10.1007/s00037-019-00186-3.

\bibitem[51]{published-51}
M. Kumar, C. Ramya, R. Saptharishi, and A. Tengse.


If VNP Is Hard, Then so Are Equations for It.


In 39th Int. Symp. on Th. Asp. of Comp. Sc. (STACS 2022),
volume 219 of Leibniz International Proceedings in Informatics
(LIPIcs), pages 44:1–44:13, Dagstuhl, Germany, 2022. Schloss Dagstuhl –
Leibniz-Zentrum für Informatik.


doi:10.4230/LIPIcs.STACS.2022.44.

\bibitem[52]{published-52}
M. Kumar and B. L. Volk.


A lower bound on determinantal complexity.


Comput. Complex., 31(2):12, 2022.


doi:10.1007/S00037-022-00228-3.

\bibitem[53]{published-53}
J. M. Landsberg, L. Manivel, and N. Ressayre.


Hypersurfaces with degenerate duals and the Geometric Complexity
Theory program.


Comment. Math. Helv., 88(2):469–484, 2013.


doi:10.4171/CMH/292.

\bibitem[54]{published-54}
J. M. Landsberg and G. Ottaviani.


New lower bounds for the border rank of matrix multiplication.


Th. of Comp., 11(11):285–298, 2015.


doi:10.4086/toc.2015.v011a011.

\bibitem[55]{published-55}
N. Limaye, Srinivasan S., and S. Tavenas.


Superpolynomial lower bounds against low-depth algebraic circuits.


In 2021 IEEE 62nd Annual Symposium on Foundations of Computer
Science (FOCS), pages 804–814. IEEE, 2021.


doi:10.1109/FOCS52979.2021.00083.

\bibitem[56]{published-56}
I. G. Macdonald.


Symmetric functions and Hall polynomials.


Oxford Mathematical Monographs. The Clarendon Press Oxford
University Press, New York, second edition, 1995.

\bibitem[57]{published-57}
M. Mahajan.


Algebraic complexity classes.


Perspectives in Computational Complexity: The Somenath Biswas
Anniversary Volume, pages 51–75, 2014.


doi:10.1007/978-3-319-05446-9\_4.

\bibitem[58]{published-58}
A. I. Molev.


Gelfand-Tsetlin bases for classical Lie algebras.


In Handbook of algebra. Vol. 4, volume 4 of Handb.
Algebr., pages 109–170. Elsevier/North-Holland, Amsterdam, 2006.


doi:10.1016/S1570-7954(06)80006-9.

\bibitem[59]{published-59}
K. D. Mulmuley.


The GCT program toward the P vs. NP problem.


Communications of the ACM, 55(6):98–107, 2012.


doi:10.1145/2184319.2184341.

\bibitem[60]{published-60}
K. D. Mulmuley and M. Sohoni.


Geometric Complexity Theory I: An Approach to the P vs. NP and
Related Problems.


SIAM J. Comput., 31(2):496–526, 2001.


doi:10.1137/S009753970038715X.

\bibitem[61]{published-61}
K. D. Mulmuley and M. Sohoni.


Geometric complexity theory II: towards explicit obstructions for
embeddings among class varieties.


SIAM Journal on Computing, 38(3):1175–1206, 2008.


doi:10.1137/080718115.

\bibitem[62]{published-62}
N. Nisan and A. Wigderson.


Lower bounds on arithmetic circuits via partial derivatives.


Computational complexity, 6:217–234, 1996.


doi:10.1007/BF01294256.

\bibitem[63]{published-63}
L. Oeding and S.V Sam.


Equations for the fifth secant variety of Segre products of
projective spaces.


Exp. Math., 25(1):94–99, 2016.


doi:10.1080/10586458.2015.1037872.

\bibitem[64]{published-64}
A. M. Perelomov and V. S. Popov.


Casimir Operators for Semisimple Lie Groups.


Mathematics of the USSR-Izvestiya, 2(6):1313, December 1968.


doi:10.1070/IM1968v002n06ABEH000731.

\bibitem[65]{published-65}
G. Racah.


Sulla caratterizzazione delle rappresentazioni irriducibili dei
gruppi semisemplici di Lie.


Atti Accad. Naz. Lincei Rend. Cl. Sci. Fis. Mat. Nat. (8),
8:108–112, 1950.

\bibitem[66]{published-66}
A. A. Razborov and S. Rudich.


Natural proofs.


In Proc. of the 26th Ann. ACM Symp. Th. of Comp., pages
204–213, 1994.

\bibitem[67]{published-67}
K. W. Regan.


Understanding the Mulmuley–Sohoni approach to P vs. NP.


Bulletin of the EATCS, 78:86–99, 2002.

\bibitem[68]{published-68}
R. Saptharishi.


A survey of lower bounds in arithmetic circuit complexity.


https://github.com/dasarpmar/lowerbounds-survey/, 2021.


Version 9.0.3.

\bibitem[69]{published-69}
V. Strassen.


Rank and optimal computation of generic tensors.


Lin. Alg. Appl., 52/53:645–685, 1983.


doi:10.1016/0024-3795(83)80041-X.

\bibitem[70]{published-70}
J. J. Sylvester.


On the principles of the calculus of forms.


Cambridge and Dublin Math. J., 7:52–97, 1852.

\bibitem[71]{published-71}
L. G. Valiant.


Completeness classes in algebra.


In Proceedings of the 11th ACM Symp. on Th. of Comp., STOC
’79, pages 249–261, New York, 1979. ACM.


doi:10.1145/800135.804419.

\bibitem[72]{published-72}
D. P. Žhelobenko.


The classical groups. Spectral analysis of their finite-dimensional
representations.


Uspehi Mat. Nauk, 17(1(103)):27–120, 1962.


doi:10.1070/RM1962v017n01ABEH001123.

\end{thebibliography}
